\section{Wstęp}
Dobór architektury sieci neuronowej
i parametrów uczenia jest procesem,
wymagającym doświadczenia.
Natomiast nawet posiadając
doświadczenie nie jest to pewne,
że dany model będzie odpowiednio działał ani,
że jest najlepszym z możliwych.
Konsekwencją tego jest potrzeba licznych
eksperymentów z różnymi modelami,
by zapewniać coraz lepiej działające modele.
Trzeba testować wiele konfiguracji warstw,
liczby neuronów, funkcji aktywacji
czy parametrów uczenia,
co oznacza dużą przestrzeń możliwych rozwiązań.
Do tego najlepiej testować
daną konfigurację kilka razy,
ponieważ są losowane startowe wagi,
które mogą wpływać znacząco
na końcowy wygląd sieci po uczeniu.

\noindent\\
Potrzeba ograniczenia ręcznego
zmieniania parametrów i testowania,
kieruje do automatyzacji tego procesu.
Do tego zostanie wykorzystany algorytm genetyczny,
ponieważ jest dobrze przystosowany
do przeszukiwania dużych przestrzeni rozwiązań.

\subsection*{Cel}
Celem pracy jest opracowanie biblioteki
w języku Python, która umożliwi automatyczne
optymalizowanie sieci neuronowej
z pomocą algorytmu genetycznemu
pod względem architektury sieci
i jej parametrów uczenia
na podstawie danych treningowych
i walidacyjnych, i testowych.

% \noindent\\
% Algorytmy genetyczne, pozwalają efektywnie przeszukiwać przestrzeń możliwych modeli
% dzieki budowaniu ich na podstawie "genów".
% Dzięki tej charakterystyce powinien dobrze sprawdzić się w~dobieraniu odpowiedniej architektury i parametrów.

% \noindent\\
% Celem pracy jest zbudowanie biblioteki,
% która automatyzuje proces dobierania przez wspomniany algorytm genetyczny.

\subsection*{Cele szczegółowe}
Pierwszym kamieniem milowym jest opracowanie struktury
kodu genetycznego umożliwiającego generowanie z nich,
w spójny sposób, architektury sieci neuronowej
oraz parametrów uczenia.

\noindent\\
Kolejnym jest zawarcie w kodzie nauczania
sieci neuronowych, które będzie względnie niezależne
od wybranego przez użytkownika platformy
(tensorflow, pytorch).

\noindent\\
Ostatnim krokiem jest raportowanie wszystkich
prób do pliku wraz z ich nauczonymi wagami.

\subsection*{Zakres}
Program powininen optymalizować
architekturę sieci neuronowej
wraz z jej parametrami uczenia
dla danych podanych przez użytkownika.
Dane będą przyjmowane w postaci
tabel pakietu \texttt{numpy}.

% is in theory?
% \noindent\\
% Przez architekturę sieci, w niniejszym dokumencie,
% rozumie się liczbę warstw,
% liczbę neuronów w~każdej z warstw
% oraz funkcje aktywacji dla każdej z warstw.
% A przez parametry uczenia
% rozumie się rodzaj optymalizatora,
% wskaźnik uczenia się oraz liczbę epok.

\noindent\\
Optymalizacja powinna się odbywać dla konkretnych
danych dostarczonych przez użytkownika,
za pomocą algorytmu genetycznego.
Implementacja algorytmu genetycznego będzie obejmowała:
osobników, selekcję, krzyżowanie i mutację.

\noindent\\
Każdy osobnik będzie przechowywał
jednoznaczny kod genetyczny, to znaczy,
że dany kod genetyczny za każdym razem
doprowadzi do takiej samej architektury
wraz z takimi samymi parametrami uczenia.

\noindent\\
Zostaną przeprowadzone eksperymenty,
gdzie będzie spradzane
czy automatycznie dobrane sieci
są lepszej lub takiej samej jakości
jak te dobrane manualnie.
Oraz zostaną określone domyślne parametry
dla algorytmu genetycznego,
które będą możliwie najkorzystniejsze
dla optymalizacji architektur sieci neuronowych.

\subsection*{Ograniczenia}
Program nie operuje na wszystkich
typach sieci neuronowych,
tworzy tylko sieci MLP,
więc po użyciu stworzonego
w niniejszej pracy programu,
można dalej zoptymalizować sieć.
Do jeszcze lepszego działania
mogą się przyczynić:
połączenia przechodzące na inne warstwy
niż tylko następna w kolejności
i/lub mniejsza ilość połączeń między neuronami
(czyli gdy przestaną być w pełni połączone).

\noindent\\
Uczenie sieci neuronowej jest operacją zużywającą
dużą ilość mocy obliczeniowej.
Dodatkowo ponieważ proces uczenia sieci zależy
od pierwotnie wylosowanych wag,
to dla uzyskania większej pewności
co do uzyskanego wyniku należy
proces uczenia wykonać kilkakrotnie.
Powoduje to, że działanie programu
jest znacząco uzależnione
od mocy obliczeniowej urządzenia
wykonującego operacje
i proces optymalizacji
poprzez algorytm genetyczny
jest stosunkowo długi.

\subsection*{Teza}
W niniejszej pracy przyjęto tezę,
że wykorzystanie algorytmu genetycznego
do optymalizacji architektury sieci neuronowej
i jej parametrów uczenia umożliwia uzyskanie modeli
o jakości porównywalnej lub lepszej
niż rozwiązania dobierane manualnie
przy ograniczeniu nakładu pracy użytkownika.

\subsection*{Struktura pracy}
Tu jeszcze nic nie ma,
ponieważ będą na bieżąco dodawane
krótkie opisy do danego
rozdziału po jego skończeniu.

\noindent\\
Aktualnie by wypełnić przestrzeń:\\
\lipsum[1][1-10]
